\scanpage{46}
\mspage{2}
\uline{Objets} : Triples $(X,\widetilde\Sigma,
\deleted{\unsure{S}}\inserted{R})$, où

$X$ surface à bord compacte orientée

$\widetilde\Sigma$ sous-variété compacte de $X$
\deleted{\unsure{close}}, $\widetilde\Sigma\supset\partial X$

$\deleted{\unsure{S}}R\subset\Sigma\deleted{\unsure{\widetilde\Sigma}}$,
partie finie

\margin{$X,\unsure{\partial X},\Sigma,S$}
\margin{$\widetilde\Sigma=\partial X\amalg\Sigma$

On peut considérer que $R$ rigidifie $\widetilde\Sigma$ :
que $R$ ``découpe'' $\widetilde\Sigma$.}

\uline{flèches} : classes d'isotopie d'homéom. conservant l'orientation

\deleted{Variante \unsure{de}}
Les composantes connexes de $\widetilde\Sigma\setminus\partial X
\inserted{=\Sigma}$ sont les \uline{cercles de découpage},
\inserted{celles de $\partial X$ sont les cercles de raccord},
les éléments de $\deleted{\unsure{S}}R$ sont les
\deleted{\unsure{pts de marquage}} \uline{sommets de
\deleted{marquage}\inserted{rigidification}},
\deleted{(ou de rigidification)},
\deleted{les composantes connexes de $\partial X$ sont les ``trous''.}
Les composantes connexes de $X$ sont les
$\deleted{\text{connexes}}$\inserted{\uline{composants}}
de l'objet envisagé.

Variante $S_*\mathcal TD_{\mathrm{top}}^\pm$, avec
\diagram{\begin{tikzcd}[ampersand replacement=\&]
S_*\mathcal TD_{\mathrm{top}}\arrow[r,phantom,"\subset"]\&
S_*\mathcal TD_{\mathrm{top}}^\pm\arrow[r]\&{[\pm1]}
\end{tikzcd}}
comme tantôt.

On a des sous-groupoïdes pleins
\diagram{\(\begin{array}{ccl}
R=\varnothing,\ \Sigma=\unsure{\varnothing}&
\mathcal T_{\mathrm{top}}\quad\text{et}\quad\mathcal T^\pm_{\mathrm{top}}&\\
\deleted{\unsure{R=\widetilde\Sigma=\varnothing}}&&\\
\Sigma=\deleted{\unsure{\partial X}}\varnothing&
S_*\mathcal T_{\mathrm{top}}\quad\text{et}\quad S_*\mathcal T^\pm_{\mathrm{top}}&
\begin{gathered}
\text{(groupoïdes de Teichmüller \uline{spéciaux}}\\
\text{\inserted{ou \uline{rigidifiés}} (topologiques orientés))}
\end{gathered}\\
R\text{ quelconque}&&\\
R=\varnothing&
\mathcal TD_{\mathrm{top}}\quad\text{et}\quad\mathcal T^\pm D_{\mathrm{top}}&
\end{array}\)}
(Cas pas intéressant je crois, sauf pour $\Sigma=\unsure{\varnothing}$
également)

\margin{MD = ``modulaire'' ! Mumford--Deligne}

\deleted{\unsure{Dé}} \uline{Détermination} du
$\pi_0(S_*\mathcal TD_{\mathrm{top}})
 \xrightarrow{\sim}\pi_0((S_*\mathcal TD)^\pm_{\mathrm{top}})$.

On introduit les \uline{graphes} \deleted{\unsure{MD}} bipondérés
\mspage{45}
